\documentclass[10pt,a4paper]{amsart}
\usepackage[margin=19mm]{geometry}
\usepackage{amsmath,amssymb,mathtools}
\usepackage[scr=rsfs,cal=euler]{mathalfa}
\usepackage{xfrac}
\usepackage[unicode,hidelinks,bookmarks=false]{hyperref}
\newcommand{\ie}{i.e.\ }
\newcommand{\cf}{cf.\ }
\newcommand{\resp}{respectively\ }
\newcommand{\Subsection}{\subsection}
\providecommand{\nobreakdash}{\nobreak\hbox{-}}
\setlength{\emergencystretch}{2em}
\multlinegap0pt
\newcommand{\F}{\mathbb F}
\usepackage{amssymb}
\usepackage{mathtools}
\usepackage{xcolor}
\newtheorem{lemma}{Lemma}
\newtheorem{proposition}{Proposition}
\newtheorem{theorem}{Theorem}
\newcommand{\Gr}{\operatorname{Gr}}
\newcommand{\supp}{\operatorname{supp}}
\title{The Structure of Cycles in Projective Geometry over $\F_q$}
\author{Ran J. Tessler, Elad Tzalik}
\date{Source: arXiv:2608.07683v1 (CC BY 4.0)}
\begin{document}
\maketitle

\noindent\emph{Source and licence.} The Structure of Cycles in Projective Geometry over $\F_q$. Version arXiv:2608.07683v1 (CC BY 4.0), distributed by arXiv under the Creative Commons Attribution 4.0 International licence, http://creativecommons.org/licenses/by/4.0/, as recorded in the arXiv licence metadata for this exact version and shown on the abstract page https://arxiv.org/abs/2608.07683v1. Full licence notice: \texttt{licenses/arxiv-2608.07683-CC-BY-4.0.txt}.\par\smallskip

\noindent\textit{Editorial scope.} Counted unit: the article's theorem thm:grassmann-coboundary-expansion (source0.tex lines 951--965). The article's complete original proof of it is delivered with it (source0.tex lines 967--1097, one \texttt{proof} environment of 131 lines). No mathematics is written, repaired or re-derived: the statement and the proof are the article's own lines, in the article's own notation, and no retained line is substituted or re-flowed.  Retained uncounted as the source-local context the proof needs: lines 780--787; lines 808--818; lines 878--891; lines 917--934.  Nothing was omitted from any retained span, so the package carries no omission record.  No retained line contains a citation, so the wrapper carries no bibliography and no bibliography entry.  No retained line contains a figure environment, an image-inclusion command or drawing code, so no image is delivered and the package declares no asset dependency.  Editorial formatting only, in addition: an AMS \texttt{amsart} preamble that restates the article's own declarations and macros used by the retained lines, namely the environments \texttt{lemma}, \texttt{proposition}, \texttt{theorem} on separate counters, together with the standard packages those lines need.  The retained environments print in this document as Lemma 1, Proposition 1, Theorem 1 in order of appearance, whereas the article prints the same items with the numbering of its own full document.  The complete native source of the article as distributed in the arXiv e-print for this exact version is bundled unchanged as \texttt{sources/207101/source0.tex}.
\begin{lemma}[Equivariance]\label{lem:cone-equivariance}
Let $g\in\operatorname{GL}(V)$ and write
$g\mathbf b=(gb_1,\ldots,gb_n)$. For every $j\ge0$ with
$2j+1\le n$ and every $x\in C_j(V)$,
\[
 c_{g\mathbf b,j}(gx)=g\,c_{\mathbf b,j}(x).
\]
\end{lemma}

\begin{lemma}[Orthogonal-complement duality]
\label{lem:orthogonal-complement-duality}
For every $0\le j\le n$, the map $x\mapsto x^\perp$ is a
support-preserving isomorphism from $C_j(V)$ to $C^{n-j}(V)$ and
satisfies
\[
 (\partial_jx)^\perp=d_{n-j}(x^\perp).
\]
Consequently,
$H_j(V)\cong H^{n-j}(V)$.
\end{lemma}

\begin{lemma}[Cochain contraction]\label{lem:cochain-contraction}
If $n\ge1$, then
\[
 \iota_{\mathbf b,0}d_0=I_{C^0(V)}.
\]
For every $j\ge1$ with $2j+1\le n$,
\[
 \iota_{\mathbf b,j}d_j
 +
 d_{j-1}\iota_{\mathbf b,j-1}
 =
 I_{C^j(V)}.
\]
\end{lemma}

\begin{proposition}[Cone size]\label{prop:cone-size}
For every ordered basis $\mathbf b$, every $j$ with $2j+1\le n$,
and every $W\in\Gr_j(V)$,
\[
 |c_{\mathbf b,j}[W]|\le\kappa_j.
\]
Moreover,
\[
 \kappa_j
 =
 [j+1]_q[j]_q!^2
 \left(
 1+\sum_{r=1}^j\frac{1}{[r]_q!^2}
 \right).
\]
In particular, $\kappa_j$ bounds the cone sizes in every degree at most
$j$.
\end{proposition}

\begin{theorem}[Coboundary expansion]
\label{thm:grassmann-coboundary-expansion}
Let $1\le k<n/2$. Suppose that
$|c_{\mathbf b,j}[W]|\le\Lambda$ for every ordered basis $\mathbf b$,
every $0\le j\le k$, and every $W\in\Gr_j(V)$. Then
\[
 h_{k,n}
 \ge
 \frac{[n-k]_q}{[k+1]_q}\frac1\Lambda,
 \qquad
 h_{n-k,n}\ge\frac1\Lambda.
\]
In particular, one may take $\Lambda=\kappa_k$ from
Proposition~\ref{prop:cone-size}.
\end{theorem}

\begin{proof}
Let $\mathcal B(V)$ be the set of ordered bases of $V$. Since $d_k$
vanishes on $B^k(V)$, it is constant on every coset of $B^k(V)$.
Fix a nonzero coset in $C^k(V)/B^k(V)$, and choose a representative
$\alpha$ of minimum support. Thus
\[
 |\alpha|
 =
 \operatorname{dist}\bigl(\alpha,B^k(V)\bigr).
\]
For every $\mathbf b\in\mathcal B(V)$,
Lemma~\ref{lem:cochain-contraction} gives
\[
 \alpha-d_{k-1}\iota_{\mathbf b,k-1}\alpha
 =
 \iota_{\mathbf b,k}d_k\alpha.
\]
The left-hand side represents the same coset as $\alpha$. Hence the
minimality of $\alpha$ implies
\[
 |\alpha|\le|\iota_{\mathbf b,k}d_k\alpha|.
\]

We next relate the support on the right to the cones. For every
$W\in\Gr_k(V)$,
\[
 (\iota_{\mathbf b,k}d_k\alpha)([W])
 =
 (d_k\alpha)(c_{\mathbf b,k}[W]).
\]
If this value is nonzero, then some member of
$\supp(c_{\mathbf b,k}[W])$ has a nonzero coefficient in
$d_k\alpha$. Possible cancellations can only make the support of the
contracted cochain smaller. Therefore
\[
 |\iota_{\mathbf b,k}d_k\alpha|
 \le
 \sum_{W\in\Gr_k(V)}
 \mathbf 1_{\{
 \supp(c_{\mathbf b,k}[W])
 \cap\supp(d_k\alpha)\ne\varnothing\}}.
\]
Summing over all $\mathbf b \in \mathcal B(V) $
\[
 |\mathcal B(V)|\,|\alpha|
 \le
 \sum_{\mathbf b\in\mathcal B(V)}
 \sum_{W\in\Gr_k(V)}
 \mathbf 1_{\{
 \supp(c_{\mathbf b,k}[W])
 \cap\supp(d_k\alpha)\ne\varnothing\}}.
\]

For $U\in\Gr_{k+1}(V)$, let
\[
 N(U)
 =
 \bigl|\{(\mathbf b,W)\in\mathcal B(V)\times\Gr_k(V):
 U\in\supp(c_{\mathbf b,k}[W])\}\bigr|.
\]
Lemma~\ref{lem:cone-equivariance} and the transitive action of
$\operatorname{GL}(V)$ on $\Gr_{k+1}(V)$ show that $N(U)$ is
independent of $U$; denote its common value by $N$. Double counting
the intersections with $\supp(d_k\alpha)$ gives
\begin{align*}
 &\sum_{\mathbf b\in\mathcal B(V)}
 \sum_{W\in\Gr_k(V)}
 \mathbf 1_{\{
 \supp(c_{\mathbf b,k}[W])
 \cap\supp(d_k\alpha)\ne\varnothing\}}\\
 &\qquad\le
 \sum_{U\in\supp(d_k\alpha)}N(U)
 =
 N|d_k\alpha|.
\end{align*}
On the other hand, counting all incidences between cones and
$(k+1)$-spaces yields
\begin{align*}
 |\Gr_{k+1}(V)|N
 &=
 \sum_{U\in\Gr_{k+1}(V)}N(U)\\
 &=
 \sum_{\mathbf b\in\mathcal B(V)}
 \sum_{W\in\Gr_k(V)}
 |c_{\mathbf b,k}[W]|\\
 &\le
 |\mathcal B(V)|\,|\Gr_k(V)|\,\Lambda.
\end{align*}
Combining these estimates and cancelling $|\mathcal B(V)|$ gives
\[
 |\alpha|
 \le
 \frac{|\Gr_k(V)|}{|\Gr_{k+1}(V)|}
 \Lambda\,|d_k\alpha|.
\]
Since
$|\Gr_{k+1}(V)|/|\Gr_k(V)|=[n-k]_q/[k+1]_q$, this proves the first
bound.

For the second bound, fix $\alpha\in C^{n-k}(V)$ and write
$\alpha=\beta^\perp$ for the unique $\beta\in C_k(V)$. Put
$z=\partial\beta\in Z_{k-1}(V)$ and fix an ordered basis $\mathbf b$.
Since $z$ is a cycle, the cone identity gives
$\partial c_{\mathbf b,k-1}(z)=z$, so
$c_{\mathbf b,k-1}(z)$ is a filling of $z$. Applying the cone identity
to $\beta$ also gives
\[
 \beta-c_{\mathbf b,k-1}(z)
 =
 \partial c_{\mathbf b,k}(\beta).
\]
Applying orthogonal complementation and
Lemma~\ref{lem:orthogonal-complement-duality}, we obtain
\[
 \alpha-\bigl(c_{\mathbf b,k-1}(z)\bigr)^\perp
 =
 d_{n-k-1}\bigl(\bigl(c_{\mathbf b,k}(\beta)\bigr)^\perp\bigr)
 \in B^{n-k}(V).
\]
Moreover, $z^\perp=d_{n-k}\alpha$. Since orthogonal complementation
preserves support and cone size, and the cone size is at most $\Lambda$ in degree
$k-1$,
\begin{align*}
 \operatorname{dist}\bigl(\alpha,B^{n-k}(V)\bigr)
 &\le |c_{\mathbf b,k-1}(z)|\\
 &\le
 \sum_{W\in\supp(z)}|c_{\mathbf b,k-1}[W]|\\
 &\le \Lambda|z|
 =\Lambda|d_{n-k}\alpha|.
\end{align*}
\end{proof}

\end{document}
